% ==============================================================================
% Chapter 5: Empirical Performance and Runtime Analysis
% ==============================================================================
\chapter{Empirical Performance and Runtime Analysis}
\label{chap:performance}

While static type systems (\cref{chap:type_systems}) and syntactic expressiveness (\cref{chap:expressiveness}) govern the compile-time correctness and structural clarity of formal language implementations, execution dynamics govern the responsiveness, scalability, and resource efficiency of interactive notebook workflows. In formal languages and compiler construction, workloads span a wide computational spectrum: from linear-time character-stream scanning ($O(n)$) to fixed-point dataflow equations ($\text{FIRST}$ and $\text{FOLLOW}$ sets), worst-case exponential subset constructions ($O(2^{|Q|})$), partition-refinement minimization ($O(|\Sigma|\,|Q|\log|Q|)$), cubic-time general context-free chart parsing ($O(n^3)$ in Earley's algorithm), and canonical $\text{LR}(1)/\text{LALR}(1)$ state-machine generation over large production grammars.

When these algorithms are hosted inside interactive Jupyter notebooks, total execution latency is determined not solely by raw CPU instruction throughput, but by the interaction of four distinct architectural layers:
\begin{enumerate}
  \item \textbf{Kernel Process Initialization and IPC Protocol Overhead:} The time required to spawn the language runtime, bind ZeroMQ sockets, initialize the read-eval-print loop (REPL) or compiler service, and serialize cell execution requests and rich MIME outputs (HTML, SVG, plain text).
  \item \textbf{Code Evaluation Strategy:} Bytecode interpretation in CPython 3.13, incremental bytecode compilation in the OCaml 5.4 top-level (\texttt{ocaml-jupyter}) versus native ahead-of-time compilation (\texttt{ocamlopt}), and incremental type-checking plus tiered Just-In-Time (RyuJIT) compilation via \texttt{FSharp.Compiler.Service} in \texttt{dotnet-interactive} (.NET~9 CoreCLR).
  \item \textbf{Data Structure Representation and Asymptotic Constant Factors:} The interplay between collection semantics and memory layout---specifically, C-backed open-addressing hash tables with cached hash codes in Python (\texttt{set}, \texttt{frozenset}, \texttt{dict}) versus pointer-linked persistent balanced binary search trees (AVL/weight-balanced trees) in OCaml (\texttt{Set.Make}, \texttt{FSet}) and F\# (\texttt{Set<'T>}, \texttt{Map<'K,'V>}).
  \item \textbf{Memory Allocation and Garbage Collection Dynamics:} Reference counting augmented by generational cycle detection in CPython, bump-pointer minor-heap allocation with mark-and-sweep major collection in OCaml~5, and generational Segregated Object Heap (SOH) management in the .NET~9 CoreCLR garbage collector.
\end{enumerate}

This chapter presents a rigorous empirical investigation of all four layers across Python~3.13, OCaml~5.4, and F\#~9.0 (.NET~9), combining isolated algorithmic microbenchmarks with end-to-end macrobenchmarks across the complete curriculum notebook corpus.

% ------------------------------------------------------------------------------
\section{Experimental Methodology and System Architecture}
\label{sec:perf_methodology}
% ------------------------------------------------------------------------------

To ensure statistical validity and eliminate cross-contamination with the source repositories, all measurements were conducted using an automated, isolated benchmarking harness (\texttt{benchmarks/microbenchmarks.py} and \texttt{benchmarks/notebook\_benchmarks.py}) following the guidelines for statistically rigorous runtime evaluation established by Georges et al.~\cite{georges2007statistically}.

\subsection{Hardware and Runtime Environment}

All experiments were executed on a unified Apple Silicon (\texttt{arm64} / \texttt{osx-arm64}) workstation under controlled thermal and background-load conditions. The exact language runtimes, compilers, and Jupyter kernel implementations evaluated are summarized in \cref{tab:runtime_configurations}.

\begin{table}[htbp]
  \centering
  \caption{Evaluated language runtimes, compilation models, and Jupyter kernel architectures.}
  \label{tab:runtime_configurations}
  \small
  \begin{tabularx}{\textwidth}{@{}llXl@{}}
    \toprule
    \textbf{Platform} & \textbf{Version} & \thead{Execution \& Compilation\\Model} & \textbf{Jupyter Kernel} \\
    \midrule
    \textbf{Python} & CPython 3.13.11 & Bytecode compiler + specializing adaptive interpreter (PEP 659) & \texttt{ipykernel} 7.4.0 \\
    \addlinespace[4pt]
    \textbf{OCaml (Toplevel)} & OCaml 5.4.1 & Incremental bytecode compiler + virtual machine (\texttt{ocamlrun}) & \texttt{ocaml-jupyter} \\
    \addlinespace[4pt]
    \textbf{OCaml (Native)} & OCaml 5.4.1 & Ahead-of-time native \texttt{arm64} machine code (\texttt{ocamlopt -O3}) & Standalone CLI \\
    \addlinespace[4pt]
    \textbf{F\# (.NET 9)} & SDK 9.0.122 & FCS + CoreCLR Tiered RyuJIT native engine (\texttt{arm64}) & \texttt{dotnet-interactive} \\
    \bottomrule
  \end{tabularx}
\end{table}

\subsection{Instrumentation and Measurement Metrics}

We partition our empirical study into three complementary experimental regimes:
\begin{description}
  \item[Regime I: Kernel Lifecycle (\cref{sec:kernel_latency})] Using programmatic kernel manager calls, we decompose cold-start notebook initialization into:
  \begin{equation}
    T_{\text{cold}} = T_{\text{spawn}} + T_{\text{cell}_1},
  \end{equation}
  where $T_{\text{spawn}}$ is the wall-clock interval from process creation until the ZeroMQ \texttt{iopub} and \texttt{shell} channels report readiness, and $T_{\text{cell}_1}$ is the latency required to compile and evaluate the first trivial expression (\texttt{1 + 1}). Subsequent evaluations of the same expression yield the warm inter-process communication (IPC) and incremental evaluation floor $T_{\text{warm}}$.

  \item[Regime II: Isolated Microbenchmarks (\cref{sec:microbenchmarks})] Because Jupyter notebooks interleave algorithmic computation with ZeroMQ IPC, HTML/SVG formatting, and external Graphviz (\texttt{dot}) subprocess invocations, notebook wall-clock timings alone cannot isolate pure language runtime performance. We therefore implemented four self-contained, mathematically identical formal-language workloads in Python~3.13, OCaml~5.4 (both bytecode and native \texttt{ocamlopt -O3}), and F\#~9.0, measuring pure algorithmic execution time via high-resolution monotonic clocks (\texttt{time.perf\_counter}, \texttt{Unix.gettimeofday}, \texttt{System.Diagnostics.Stopwatch}) and peak Resident Set Size ($\text{RSS}_{\max}$) via BSD/Mach \texttt{getrusage(RUSAGE\_SELF)} (\texttt{/usr/bin/time -l}). Each microbenchmark was executed across $N = 5$ independent process invocations, reporting the arithmetic mean $\mu$, sample standard deviation $\sigma$, median, and extrema.

  \item[Regime III: Notebook Macrobenchmarks (\cref{sec:macrobenchmarks})] Using a custom subclass of \texttt{nbclient.NotebookClient} (\texttt{TimedNotebookClient}), we executed 14 representative matched notebooks spanning Chapters~03 through~13 inside an ephemeral workspace copy. For each notebook $B$, we record:
  \begin{align}
    T_{\text{wall}}(B) &= t_{\text{teardown}} - t_{\text{init}}, \\
    T_{\text{cells}}(B) &= \sum_{c \in \text{CodeCells}(B)} \bigl(t_{\text{end}}(c) - t_{\text{start}}(c)\bigr), \\
    T_{\text{overhead}}(B) &= T_{\text{wall}}(B) - T_{\text{cells}}(B),
  \end{align}
  alongside the cumulative Resident Set Size ($\text{RSS}_{\text{tree}}$) of the kernel process and all child processes immediately prior to kernel shutdown.
\end{description}

% ------------------------------------------------------------------------------
\section{Kernel Initialization and IPC Protocol Latency}
\label{sec:kernel_latency}
% ------------------------------------------------------------------------------

Before a single line of user code executes in a Jupyter notebook, the host server must spawn the language kernel process, establish five ZeroMQ sockets (\texttt{shell}, \texttt{iopub}, \texttt{stdin}, \texttt{control}, and \texttt{hb}), complete the \texttt{kernel\_info\_request} handshake, and initialize the language's interactive compiler front-end. \Cref{tab:kernel_cold_start} and \cref{fig:kernel_cold_start} present the empirical latency decomposition and baseline memory footprint across the three kernels.

\begin{table}[htbp]
  \centering
  \caption{Jupyter kernel cold-start latency decomposition, warm IPC roundtrip time, and baseline idle Resident Set Size (mean $\pm$ standard deviation).}
  \label{tab:kernel_cold_start}
  \small
  \setlength{\tabcolsep}{5.2pt}
  \begin{tabular}{@{}lrrrrr@{}}
    \toprule
    \thead{Kernel Platform} &
    \thead{$T_{\text{spawn}}$\\(ms)} &
    \thead{$T_{\text{cell}_1}$\\(ms)} &
    \thead{$T_{\text{cold}}$\\(ms)} &
    \thead{$T_{\text{warm}}$\\(ms)} &
    \thead{Idle RSS\\(MB)} \\
    \midrule
    Python 3.13 (\texttt{ipykernel}) &
      $384.5 \pm 1.2$ &
      $11.2 \pm 0.9$ &
      $395.7$ &
      $4.49 \pm 0.25$ &
      $55.4$ \\
    OCaml 5.4 (\texttt{ocaml-jupyter}) &
      $504.0 \pm 95.6$ &
      $7.2 \pm 0.3$ &
      $511.1$ &
      $3.50 \pm 0.33$ &
      $27.2$ \\
    F\# 9.0 (\texttt{.net-fsharp}) &
      $863.2 \pm 162.9$ &
      $1021.2 \pm 5.2$ &
      $1884.4$ &
      $3.70 \pm 0.18$ &
      $530.5$ \\
    \bottomrule
  \end{tabular}
\end{table}

\begin{figure}[htbp]
  \centering
  \begin{tikzpicture}
    \begin{axis}[
      ybar stacked,
      bar width=24pt,
      width=0.88\textwidth,
      height=6.8cm,
      ylabel={Latency (ms)},
      symbolic x coords={Python 3.13, OCaml 5.4, F\# 9.0},
      xtick=data,
      xticklabel style={font=\small, align=center},
      ymin=0,
      ymax=2200,
      ymajorgrids=true,
      grid style={dashed,gray!40},
      legend style={at={(0.03,0.95)}, anchor=north west, font=\small},
      nodes near coords,
      every node near coord/.append style={font=\scriptsize},
      enlarge x limits=0.25
    ]
      \addplot+[fill=PythonBlue!80, draw=PythonBlue] coordinates {
        (Python 3.13, 384.5)
        (OCaml 5.4, 504.0)
        (F\# 9.0, 863.2)
      };
      \addplot+[fill=OCamlOrange!80, draw=OCamlOrange] coordinates {
        (Python 3.13, 11.2)
        (OCaml 5.4, 7.2)
        (F\# 9.0, 1021.2)
      };
      \legend{Kernel Spawn \& ZeroMQ Handshake ($T_{\text{spawn}}$), First-Cell Compiler Warmup ($T_{\text{cell}_1}$)}
    \end{axis}
  \end{tikzpicture}
  \caption{Decomposition of total cold-start latency ($T_{\text{cold}} = T_{\text{spawn}} + T_{\text{cell}_1}$) across the three Jupyter kernels.}
  \label{fig:kernel_cold_start}
\end{figure}

\subsection{Analysis of Initialization Mechanics}

The measurements in \cref{tab:kernel_cold_start} reveal fundamental architectural differences among the three interactive execution engines:

\begin{enumerate}
  \item \textbf{CPython \texttt{ipykernel} ($T_{\text{cold}} = 395.7\text{ ms}$, $\text{RSS} = 55.4\text{ MB}$):}
  Because Jupyter itself is written in Python, \texttt{ipykernel} initializes a single CPython interpreter process, loads pre-compiled \texttt{.pyc} bytecode files for \texttt{traitlets}, \texttt{pyzmq}, and \texttt{IPython}, and evaluates cell inputs directly via Python's built-in \texttt{compile()} and \texttt{eval()}/\texttt{exec()} C API. No external type-checking database or assembly metadata table is constructed, yielding a first-cell latency of only $11.2\text{ ms}$.

  \item \textbf{OCaml \texttt{ocaml-jupyter} ($T_{\text{cold}} = 511.1\text{ ms}$, $\text{RSS} = 27.2\text{ MB}$):}
  The \texttt{ocaml-jupyter} kernel is compiled to a standalone native \texttt{arm64} binary that embeds the OCaml top-level bytecode compiler (\texttt{Toploop}). Upon launch, it executes \texttt{opam env} and processes \texttt{\~/.ocamlinit} ($T_{\text{spawn}} = 504.0\text{ ms}$). Once initialized, incremental Hindley--Milner type inference and bytecode compilation of a trivial cell require only $7.2\text{ ms}$---faster even than CPython's first cell---while consuming half the memory ($27.2\text{ MB}$ RSS). Moreover, its warm cell roundtrip time ($3.50\text{ ms}$) is the lowest of all three platforms.

  \item \textbf{F\# \texttt{dotnet-interactive} ($T_{\text{cold}} = 1884.4\text{ ms}$, $\text{RSS} = 530.5\text{ MB}$):}
  The \texttt{.net-fsharp} kernel exhibits a $4.76\times$ higher cold-start latency and a $9.58\times$ larger baseline memory footprint than Python. This overhead stems from the heavyweight architecture of \texttt{Microsoft.dotnet-interactive}:
  \begin{itemize}
    \item During $T_{\text{spawn}}$ ($863.2\text{ ms}$), the CoreCLR runtime initializes its multi-language polyglot host (supporting C\#, F\#, PowerShell, HTML, and Mermaid sub-kernels simultaneously), starts the ASP.NET/ZeroMQ message routing pipeline, and JIT-compiles the host infrastructure.
    \item During $T_{\text{cell}_1}$ ($1021.2\text{ ms}$), the F\# sub-kernel lazily initializes \texttt{FSharp.\allowbreak Compiler.\allowbreak Service}, resolves BCL assemblies, constructs global symbol tables, and emits a dynamic CLI assembly via \texttt{Reflection.Emit}.
    \item Crucially, however, \emph{once this one-time $1.88\text{ s}$ initialization cost is amortized}, subsequent warm cells execute in $T_{\text{warm}} = 3.70 \pm 0.18\text{ ms}$, outperforming Python's warm IPC roundtrip ($4.49\text{ ms}$).
  \end{itemize}
\end{enumerate}

\begin{finding}[Amortized vs.\ Cold Kernel Latency]
\texttt{dotnet-interactive} incurs a fixed $\approx 1.88\text{ s}$ and $\approx 530\text{ MB}$ cold-start penalty due to .NET BCL metadata loading and \texttt{FSharp.Compiler.Service} initialization. In interactive sessions where a kernel remains alive for hours, this one-time cost is negligible ($3.70\text{ ms}$ per warm cell); in batch validation pipelines that spawn a fresh kernel per notebook, this fixed cost dominates total runtime for short notebooks.
\end{finding}

% ------------------------------------------------------------------------------
\section{Algorithmic Microbenchmarks: Pure Compute Throughput}
\label{sec:microbenchmarks}
% ------------------------------------------------------------------------------

To evaluate the intrinsic efficiency of the three languages on formal-language workloads without Jupyter IPC distortion, we benchmarked four canonical algorithms that exercise distinct runtime subsystems:
\begin{enumerate}
  \item \textbf{Workload M1 (Lexical Scanning):} Tokenizing a $900\text{ KB}$ synthetic source stream containing exactly $300{,}000$ non-whitespace lexical tokens (identifiers, integer literals, arithmetic and relational operators). This tests tight character-loop iteration, string indexing, and regular-expression/DFA scanner dispatch.
  \item \textbf{Workload M2 ($\text{NFA} \to \text{DFA}$ Powerset Construction):} Determinizing the classic parameterized family of non-deterministic automata $N_k$ recognizing $(0 \mid 1)^* 1 (0 \mid 1)^{k-1}$ for $k = 11$, which produces exactly $2^{k+1} = 4{,}096$ reachable subset states ($Q_{\text{DFA}} \subseteq \mathcal{P}(Q_{\text{NFA}})$) and $8{,}192$ subset transitions. This stresses immutable set construction, set union, and set-of-sets membership lookups (\texttt{Set<Set<int>>} vs.\ \texttt{frozenset}).
  \item \textbf{Workload M3 (DFA Partition-Refinement Minimization):} Computing the Myhill--Nerode equivalence classes of a dense $1{,}024$-state deterministic automaton down to its $5$ minimal equivalence blocks via fixed-point signature partitioning. This stresses tuple hashing/comparison, map updates, and array indexing.
  \item \textbf{Workload M4 (Symbolic AST Differentiation \& Algebraic Rewriting):} Computing the second symbolic derivative $\frac{d^2}{dx^2} E(x)$ interleaved with algebraic simplification rules ($0 + e \to e$, $0 \cdot e \to 0$, $1 \cdot e \to e$, constant folding) over a deeply nested polynomial expression tree of depth~13, yielding a simplified AST of $3{,}741$ nodes after allocating and discarding tens of thousands of intermediate AST nodes. This stresses algebraic data type allocation, nested structural pattern matching, and young-generation garbage collection.
\end{enumerate}

\Cref{tab:microbenchmarks} and \cref{fig:microbenchmarks} report the empirical results across all four execution platforms. Across all four platforms, the computed mathematical invariants were verified to be identical ($300{,}000$ tokens, $4{,}096$ DFA states, $5$ minimal partition blocks, and $3{,}741$ simplified AST nodes).

\begin{table}[htbp]
  \centering
  \caption{Isolated algorithmic microbenchmark execution times (ms, mean $\pm$ standard deviation over $N = 5$ runs) and peak process Resident Set Size (MB). Best performance in each column is highlighted in \textbf{bold}.}
  \label{tab:microbenchmarks}
  \footnotesize
  \setlength{\tabcolsep}{3.8pt}
  \begin{tabular}{@{}lrrrr@{}}
    \toprule
    \thead{Workload / Metric} &
    \thead{Python 3.13} &
    \thead{OCaml 5.4\\(Bytecode)} &
    \thead{OCaml 5.4\\(Native)} &
    \thead{F\# 9.0\\(CoreCLR)} \\
    \midrule
    \textbf{M1: Lexical Scanning} ($300{,}000$ tok) &
      $62.068 \pm 1.678$ &
      $23.862 \pm 7.848$ &
      $3.507 \pm 0.930$ &
      $\mathbf{1.915 \pm 0.027}$ \\
    \addlinespace[3pt]
    \textbf{M2: $\text{NFA} \to \text{DFA}$ Subset} ($4{,}096$ st) &
      $\mathbf{9.147 \pm 0.074}$ &
      $44.427 \pm 3.838$ &
      $10.719 \pm 2.164$ &
      $58.460 \pm 0.493$ \\
    \addlinespace[3pt]
    \textbf{M3: DFA Minimization} ($1{,}024 \to 5$ st) &
      $1.121 \pm 0.045$ &
      $1.854 \pm 0.086$ &
      $\mathbf{0.401 \pm 0.063}$ &
      $3.345 \pm 0.031$ \\
    \addlinespace[3pt]
    \textbf{M4: Symbolic AST Diff} ($3{,}741$ nodes) &
      $1.420 \pm 0.024$ &
      $0.093 \pm 0.003$ &
      $\mathbf{0.015 \pm 0.002}$ &
      $0.111 \pm 0.003$ \\
    \midrule
    \textbf{Peak Process RSS (MB)} &
      $21.79 \pm 0.04$ &
      $9.05 \pm 0.02$ &
      $\mathbf{8.77 \pm 0.00}$ &
      $411.41 \pm 0.74$ \\
    \bottomrule
  \end{tabular}
\end{table}

\begin{figure}[htbp]
  \centering
  \begin{tikzpicture}
    \begin{axis}[
      ybar,
      ymode=log,
      log origin=infty,
      bar width=9pt,
      width=0.95\textwidth,
      height=7.2cm,
      ylabel={Execution Time (ms, log scale)},
      symbolic x coords={M1: Lexing, M2: NFA to DFA, M3: DFA Min, M4: AST Diff},
      xtick=data,
      xticklabel style={font=\footnotesize},
      ymin=0.008,
      ymax=150,
      ymajorgrids=true,
      grid style={dashed,gray!40},
      legend style={at={(0.5,1.04)}, anchor=south, legend columns=4, font=\scriptsize},
      enlarge x limits=0.18
    ]
      \addplot+[fill=PythonBlue!85, draw=PythonBlue] coordinates {
        (M1: Lexing, 62.068)
        (M2: NFA to DFA, 9.147)
        (M3: DFA Min, 1.121)
        (M4: AST Diff, 1.420)
      };
      \addplot+[fill=OCamlOrange!50, draw=OCamlOrange] coordinates {
        (M1: Lexing, 23.862)
        (M2: NFA to DFA, 44.427)
        (M3: DFA Min, 1.854)
        (M4: AST Diff, 0.093)
      };
      \addplot+[fill=OCamlOrange!95, draw=OCamlOrange!80!black] coordinates {
        (M1: Lexing, 3.507)
        (M2: NFA to DFA, 10.719)
        (M3: DFA Min, 0.401)
        (M4: AST Diff, 0.015)
      };
      \addplot+[fill=FSharpPurple!85, draw=FSharpPurple] coordinates {
        (M1: Lexing, 1.915)
        (M2: NFA to DFA, 58.460)
        (M3: DFA Min, 3.345)
        (M4: AST Diff, 0.111)
      };
      \legend{Python 3.13, OCaml 5.4 (Byte), OCaml 5.4 (Native), F\# 9.0 (.NET 9)}
    \end{axis}
  \end{tikzpicture}
  \caption{Logarithmic comparison of algorithmic microbenchmark execution times across Python~3.13, OCaml~5.4 (Bytecode and Native), and F\#~9.0 (.NET~9 CoreCLR).}
  \label{fig:microbenchmarks}
\end{figure}

\subsection{Dissecting the Microbenchmark Anomalies and Runtime Physics}

The microbenchmark results in \cref{tab:microbenchmarks} expose three profound, non-obvious performance phenomena that directly impact formal-language algorithm design:

\subsubsection{1. Why F\# RyuJIT and OCaml Native Dominate Lexical Scanning (M1) by $18\times$--$32\times$}
In Workload~M1, F\#~9.0 completes scanning $300{,}000$ tokens in $1.915\text{ ms}$ (a throughput of $156.7 \times 10^6\text{ tokens/s}$), followed closely by OCaml~5.4 Native at $3.507\text{ ms}$ ($85.5 \times 10^6\text{ tokens/s}$) and OCaml Bytecode at $23.862\text{ ms}$. Python~3.13, despite delegating regular-expression matching to CPython's C-level \texttt{sre} engine via \texttt{re.finditer}, requires $62.068\text{ ms}$ ($32.4\times$ slower than F\#).

The bottleneck in Python is object allocation and boxing at the C/Python boundary: for every single token matched by \texttt{finditer}, CPython allocates a new heap-allocated \texttt{re.Match} object, increments reference counts, and dispatches dynamic attribute lookup (\texttt{mo.lastgroup}). In contrast, F\#'s tail-recursive scanner compiles under .NET~9's RyuJIT into a single tight \texttt{arm64} machine-code loop with register-allocated loop counters and range-check elimination, performing zero heap allocations across the entire $900\text{ KB}$ buffer.

\subsubsection{2. The Set-of-Sets Paradox in Powerset Construction (M2): Hash Tables vs.\ Persistent Balanced Trees}
Remarkably, in Workload~M2 ($\text{NFA} \to \text{DFA}$ subset construction), \textbf{Python~3.13 ($9.147\text{ ms}$) outperforms both OCaml Bytecode ($44.427\text{ ms}$) and F\#~9.0 ($58.460\text{ ms}$), and slightly edges out even OCaml Native ($10.719\text{ ms}$)!}

To understand this counter-intuitive result, consider the data-structural operations performed during subset construction for an NFA with $|Q| = k$ states and $2^k$ reachable subset states $S \subseteq Q$:
\begin{itemize}
  \item \textbf{Python's \texttt{frozenset} (C-level Open-Addressing Hash Table with Cached Hash):}
  In CPython, a \texttt{frozenset} stores an internal C field \texttt{hash} (of type \texttt{Py\_hash\_t}) that is computed once in $O(|S|)$ time and cached permanently on the struct. When checking whether a newly constructed subset state $S' \in \text{visited}$ (where $\text{visited}$ is a \texttt{set[frozenset[int]]}), CPython reads the cached 64-bit integer \texttt{S'->hash} in $O(1)$ time and probes the open-addressing hash table in $O(1)$ amortized time. Full set equality is invoked only upon hash collision. Consequently, inserting all $2^k$ states into $\text{visited}$ requires only $O(k \cdot 2^k)$ total elemental operations in optimized C code.

  \item \textbf{OCaml's \texttt{Set.Make(IntSet)} and F\#'s \texttt{Set<Set<int>>} (Persistent Balanced AVL Trees):}
  In both OCaml and F\#, standard immutable sets are implemented as height-balanced binary search trees (AVL trees with balance factor $\le 2$) to support purely functional persistence and structural sharing~\cite{okasaki1999purely}. Maintaining the outer set $\text{visited}$ of cardinality $M = 2^k$ requires $O(\log M) = O(k)$ tree-node comparisons per lookup or insertion. However, \emph{each comparison between two subset states $S_1, S_2 \subseteq Q$ must traverse their respective inner trees lexicographically}, costing $O(|S|) = O(k)$ steps in the worst case! Thus, checking $S' \in \text{visited}$ incurs:
  \begin{equation}
    T_{\text{lookup}}^{\text{Tree}}(M, k) = O(k \cdot \log M) = O(k^2) \quad \text{vs.} \quad T_{\text{lookup}}^{\text{Hash}}(M, k) = O(1)\text{ amortized.}
  \end{equation}
  Furthermore, while OCaml's functor \texttt{Set.Make(Int)} monomorphizes the integer comparison into a direct machine register comparison (\texttt{cmp}) yielding $10.719\text{ ms}$ in native code, F\#'s polymorphic \texttt{Set<Set<int>>} relies on .NET's generic \texttt{IComparer<'T>} interface dispatch (\texttt{LanguagePrimitives.GenericComparer}), incurring virtual interface call overhead at every node of both the outer and inner AVL trees ($58.460\text{ ms}$).
\end{itemize}

\begin{finding}[Hash Sets vs.\ Persistent Tree Sets in Automata Algorithms]
When automata algorithms construct sets of sets ($\mathcal{P}(\mathcal{P}(Q))$, such as subset states or canonical $\text{LR}(1)$ item sets) without requiring persistent historical snapshots, Python's C-level \texttt{frozenset} achieves $O(1)$ cached-hash lookups that outperform F\#'s generic AVL-tree \texttt{Set<Set<'T>>} by $6.4\times$ and OCaml's bytecode AVL-tree by $4.9\times$. Only OCaml's native functor-monomorphized \texttt{Set.Make} matches Python's speed ($10.7\text{ ms}$ vs.\ $9.1\text{ ms}$).
\end{finding}

\subsubsection{Scaling Behavior Across Exponential State Blowup ($2^8$ to $2^{13}$ States)}

To verify how this asymptotic factor scales as the deterministic state space $|Q_{\text{DFA}}| = 2^{k+1}$ grows from $256$ ($k=7$) to $8{,}192$ ($k=12$) states, we measured execution times across all six problem sizes (\cref{fig:nfa2dfa_scaling}). Across the entire 32-fold increase in state space, Python's cached-hash \texttt{frozenset} maintains the lowest absolute runtime ($19.96\text{ ms}$ at $8{,}192$ states), closely tracked by OCaml~5.4 Native ($24.58\text{ ms}$), whereas OCaml~5.4 Bytecode ($97.93\text{ ms}$) and F\#~9.0 ($142.26\text{ ms}$) diverge due to repeated $O(k^2)$ balanced-tree structural comparisons.

\begin{figure}[htbp]
  \centering
  \begin{tikzpicture}
    \begin{axis}[
      width=0.88\textwidth,
      height=6.8cm,
      xlabel={Generated DFA State Space $|Q_{\text{DFA}}| = 2^{k+1}$},
      ylabel={Execution Time (ms, log scale)},
      xmode=log,
      ymode=log,
      log basis x=2,
      xtick={256,512,1024,2048,4096,8192},
      xticklabels={$256$,$512$,$1{,}024$,$2{,}048$,$4{,}096$,$8{,}192$},
      grid=both,
      grid style={dashed,gray!35},
      legend style={at={(0.03,0.95)}, anchor=north west, font=\footnotesize},
      thick
    ]
      \addplot+[mark=*, color=PythonBlue, line width=1.2pt] coordinates {
        (256, 0.595) (512, 1.101) (1024, 2.564) (2048, 4.926) (4096, 10.089) (8192, 19.956)
      };
      \addplot+[mark=square*, color=OCamlOrange!90!black, line width=1.2pt] coordinates {
        (256, 0.623) (512, 1.396) (1024, 2.963) (2048, 6.217) (4096, 12.297) (8192, 24.584)
      };
      \addplot+[mark=triangle*, color=OCamlOrange!60, dashed, line width=1.2pt] coordinates {
        (256, 3.267) (512, 6.964) (1024, 13.611) (2048, 25.486) (4096, 47.302) (8192, 97.934)
      };
      \addplot+[mark=diamond*, color=FSharpPurple, line width=1.2pt] coordinates {
        (256, 4.229) (512, 6.540) (1024, 30.414) (2048, 32.698) (4096, 78.705) (8192, 142.262)
      };
      \legend{Python 3.13 (\texttt{frozenset}), OCaml 5.4 (Native \texttt{Set.Make}), OCaml 5.4 (Bytecode \texttt{Set.Make}), F\# 9.0 (\texttt{Set<Set<int>>})}
    \end{axis}
  \end{tikzpicture}
  \caption{Log-log scaling of $\text{NFA} \to \text{DFA}$ subset construction time as a function of reachable subset states $|Q_{\text{DFA}}| \in [256, 8{,}192]$.}
  \label{fig:nfa2dfa_scaling}
\end{figure}

\subsubsection{3. Symbolic AST Manipulation and Pattern Matching (M4): $95\times$ Advantage for ML Variants}
When the workload shifts from flat hash-set lookups to recursive Algebraic Data Type (ADT) construction and nested pattern matching (Workload~M4: computing and simplifying $\frac{d^2}{dx^2} E(x)$), the performance hierarchy reverses dramatically (\cref{lst:ast_rewrite_comparison}):
\begin{itemize}
  \item \textbf{OCaml~5.4 Native completes the entire differentiation and simplification pass in $0.015\text{ ms}$ ($15\text{ }\mu\text{s}$)}, and \textbf{OCaml~5.4 Bytecode completes it in $0.093\text{ ms}$ ($93\text{ }\mu\text{s}$)}. In OCaml, a variant constructor such as \texttt{Add(l, r)} is represented as a contiguous 3-word block allocated via a single pointer-bump instruction on the minor heap, and nested \texttt{match} expressions are compiled by the Maranget pattern-matching compiler into optimal decision trees using integer tag switches.
  \item \textbf{F\#~9.0 completes the pass in $0.111\text{ ms}$ ($111\text{ }\mu\text{s}$)}, $12.8\times$ faster than Python. In F\#, each discriminated union case compiles to a .NET subclass with a flat integer \texttt{Tag} property, enabling RyuJIT to emit fast tag-test branches and allocate short-lived AST nodes in CoreCLR's Ephemeral Gen~0 segment.
  \item \textbf{Python~3.13 requires $1.420\text{ ms}$} ($94.7\times$ slower than OCaml Native, $15.3\times$ slower than OCaml Bytecode, and $12.8\times$ slower than F\#), even when optimized with \texttt{\_\_slots\_\_} and PEP~634 \texttt{match/case}. In CPython, PEP~634 structural pattern matching executes dynamic \texttt{isinstance()} checks and attribute extractions at every AST node, while every intermediate node allocation invokes the full \texttt{PyObject} allocator and reference-count machinery.
\end{itemize}

\begin{lstlisting}[language=ML, caption={Algebraic AST simplification in OCaml/F\# (compiled jump-table dispatch over unboxed/tagged variants) versus Python PEP~634 structural pattern matching.}, label={lst:ast_rewrite_comparison}]
(* OCaml 5.4 / F# 9.0: Direct tag dispatch on algebraic variant types *)
let rec simplify e =
  match e with
  | Add (l, r) ->
      match simplify l, simplify r with
      | Num 0, r1 -> r1
      | l1, Num 0 -> l1
      | Num a, Num b -> Num (a + b)
      | l1, r1 -> Add (l1, r1)
  | Mul (l, r) ->
      match simplify l, simplify r with
      | Num 0, _ | _, Num 0 -> Num 0
      | Num 1, r1 -> r1
      | l1, Num 1 -> l1
      | Num a, Num b -> Num (a * b)
      | l1, r1 -> Mul (l1, r1)
  | _ -> e
\end{lstlisting}

% ------------------------------------------------------------------------------
\section{End-to-End Jupyter Notebook Macrobenchmarks}
\label{sec:macrobenchmarks}
% ------------------------------------------------------------------------------

While microbenchmarks isolate individual algorithmic kernels, real-world execution of the Formal Languages curriculum occurs inside Jupyter notebooks that combine module loading, parser compilation, algorithm execution, test assertions, and Graphviz SVG diagram generation. \Cref{tab:macrobenchmarks} presents the complete end-to-end execution measurements across 14 matched curriculum notebooks spanning Chapters~03 through~13, decomposing total wall-clock time ($T_{\text{wall}}$) into pure cell execution time ($T_{\text{cells}}$) and kernel startup/teardown overhead ($T_{\text{overhead}}$). \Cref{fig:macrobenchmarks_cells} visualizes the pure cell execution times across chapters.

\begin{table}[htbp]
  \centering
  \caption{End-to-end Jupyter notebook macrobenchmarks across 14 matched curriculum notebooks (steady-state warm system caches). $T_{\text{wall}}$ denotes total wall-clock time including kernel spawn; $T_{\text{cells}}$ denotes cumulative code-cell execution time; RSS denotes process-tree memory residency in MB.}
  \label{tab:macrobenchmarks}
  \scriptsize
  \setlength{\tabcolsep}{2.2pt}
  \begin{tabular}{@{}lrrrrrrrrr@{}}
    \toprule
    & \multicolumn{3}{c}{\textbf{Python 3.13 (\texttt{ipykernel})}}
    & \multicolumn{3}{c}{\textbf{OCaml 5.4 (\texttt{ocaml-jupyter})}}
    & \multicolumn{3}{c}{\textbf{F\# 9.0 (\texttt{.net-fsharp})}} \\
    \cmidrule(lr){2-4} \cmidrule(lr){5-7} \cmidrule(l){8-10}
    \thead{Curriculum Notebook} &
    \thead{$T_{\text{wall}}$\\(ms)} & \thead{$T_{\text{cells}}$\\(ms)} & \thead{RSS\\(MB)} &
    \thead{$T_{\text{wall}}$\\(ms)} & \thead{$T_{\text{cells}}$\\(ms)} & \thead{RSS\\(MB)} &
    \thead{$T_{\text{wall}}$\\(ms)} & \thead{$T_{\text{cells}}$\\(ms)} & \thead{RSS\\(MB)} \\
    \midrule
    Ch 03: Lexical Scanning &
      $578.2$ & $64.0$ & $55.5$ &
      $1{,}108.8$ & $390.6$ & $102.8$ &
      $2{,}415.4$ & $1{,}329.3$ & $575.2$ \\
    Ch 03: RegExp Tutorial &
      $618.2$ & $111.8$ & $57.1$ &
      $1{,}507.8$ & $959.2$ & $103.0$ &
      $2{,}407.7$ & $1{,}495.2$ & $596.7$ \\
    Ch 03: Exam Evaluation$^{\dagger}$ &
      $724.8$ & $229.1$ & $105.5$ &
      $1{,}108.6$ & $436.5$ & $102.8$ &
      $2{,}479.1$ & $1{,}575.7$ & $620.2$ \\
    Ch 04--05: NFA to DFA &
      $590.9$ & $56.2$ & $55.5$ &
      $874.7$ & $353.0$ & $102.6$ &
      $2{,}341.7$ & $1{,}448.3$ & $607.5$ \\
    Ch 04--05: RegExp to NFA &
      $570.9$ & $74.6$ & $56.1$ &
      $931.0$ & $389.8$ & $102.6$ &
      $2{,}441.2$ & $1{,}428.7$ & $614.3$ \\
    Ch 04--05: DFA to RegExp &
      $552.5$ & $34.0$ & $55.5$ &
      $868.6$ & $323.7$ & $102.2$ &
      $2{,}267.8$ & $1{,}372.6$ & $595.8$ \\
    Ch 04--05: DFA Minimization &
      $572.0$ & $60.8$ & $55.8$ &
      $1{,}075.1$ & $368.0$ & $104.5$ &
      $2{,}359.7$ & $1{,}497.9$ & $612.8$ \\
    Ch 04--05: Equivalence Check &
      $1{,}002.0$ & $489.4$ & $91.3$ &
      $903.2$ & $394.6$ & $102.8$ &
      $2{,}759.0$ & $1{,}669.0$ & $670.9$ \\
    Ch 06: Top-Down Parser &
      $593.9$ & $56.6$ & $55.8$ &
      $1{,}126.2$ & $463.2$ & $103.0$ &
      $2{,}353.2$ & $1{,}320.5$ & $565.1$ \\
    Ch 06: EBNF Parser &
      $551.6$ & $48.8$ & $55.5$ &
      $1{,}152.4$ & $454.4$ & $102.9$ &
      $2{,}231.0$ & $1{,}293.8$ & $563.9$ \\
    Ch 07: Differentiator &
      $583.2$ & $87.1$ & $58.7$ &
      $1{,}270.3$ & $589.4$ & $104.4$ &
      $2{,}517.2$ & $1{,}567.5$ & $612.0$ \\
    Ch 07/08: Interpreter &
      $1{,}492.7$ & $961.9$ & $92.4$ &
      $1{,}770.8$ & $1{,}058.7$ & $103.2$ &
      $2{,}932.9$ & $1{,}881.0$ & $721.3$ \\
    Ch 08/09: Earley Parser &
      $604.5$ & $108.3$ & $57.9$ &
      $1{,}188.6$ & $533.5$ & $103.3$ &
      $2{,}474.1$ & $1{,}634.3$ & $658.8$ \\
    Ch 09/10: SLR Generator$^{\ddagger}$ &
      $1{,}567.8$ & $1{,}058.9$ & $68.5$ &
      $4{,}643.6$ & $3{,}926.6$ & $242.0$ &
      $15{,}884.8$ & $14{,}842.0$ & $1{,}597.5$ \\
    Ch 10/11: Parser Conflicts &
      $1{,}243.0$ & $647.6$ & $91.4$ &
      $1{,}365.5$ & $832.0$ & $102.7$ &
      $2{,}704.0$ & $1{,}747.0$ & $607.7$ \\
    Ch 12/13: Integer-C Compiler &
      $2{,}303.4$ & $1{,}752.2$ & $92.6$ &
      $2{,}103.1$ & $1{,}429.2$ & $107.4$ &
      $2{,}907.7$ & $1{,}892.1$ & $651.7$ \\
    \midrule
    \textbf{Median (16 Notebooks)} &
      $\mathbf{599.2}$ & $\mathbf{97.7}$ & $\mathbf{57.5}$ &
      $1{,}169.4$ & $457.4$ & $102.9$ &
      $2{,}476.6$ & $1{,}532.7$ & $612.4$ \\
    \textbf{Mean (Excl.\ SLR Outlier)} &
      $\mathbf{838.7}$ & $\mathbf{318.8}$ & $\mathbf{69.1}$ &
      $1{,}223.7$ & $598.2$ & $103.3$ &
      $2{,}507.4$ & $1{,}543.3$ & $618.3$ \\
    \bottomrule
    \multicolumn{10}{@{}p{0.98\textwidth}@{}}{\scriptsize $^{\dagger}$On a cold virtualenv without a pre-built \texttt{matplotlib} font cache, Python's \texttt{03-Exam-Evaluation.ipynb} requires $10{,}170.6\text{ ms}$ ($T_{\text{cells}} = 9{,}673.9\text{ ms}$) due to macOS system font indexing. $^{\ddagger}$In the F\# edition, \texttt{SLR-Table-Generator.ipynb} additionally computes the full canonical item sets and conflict tables for the unsimplified ANSI C grammar (\texttt{c-grammar.g}) alongside \texttt{c-grammar-slr.g}.}
  \end{tabular}
\end{table}

\begin{figure}[htbp]
  \centering
  \begin{tikzpicture}
    \begin{axis}[
      xbar,
      bar width=4.5pt,
      width=0.82\textwidth,
      height=11.2cm,
      xlabel={Pure Cell Execution Time $T_{\text{cells}}$ (ms)},
      symbolic y coords={
        Ch 12/13: Integer-C Compiler,
        Ch 10/11: Parser Conflicts,
        Ch 08/09: Earley Parser,
        Ch 07/08: Interpreter,
        Ch 07: Differentiator,
        Ch 06: EBNF Parser,
        Ch 06: Top-Down Parser,
        Ch 04--05: Equivalence Check,
        Ch 04--05: DFA Minimization,
        Ch 04--05: DFA to RegExp,
        Ch 04--05: RegExp to NFA,
        Ch 04--05: NFA to DFA,
        Ch 03: Exam Evaluation,
        Ch 03: RegExp Tutorial,
        Ch 03: Lexical Scanning
      },
      ytick=data,
      xmin=0,
      xmax=2100,
      xmajorgrids=true,
      grid style={dashed,gray!40},
      legend style={at={(0.97,0.06)}, anchor=south east, font=\scriptsize},
      enlarge y limits=0.04,
      yticklabel style={font=\scriptsize}
    ]
      \addplot+[fill=PythonBlue!85, draw=PythonBlue] coordinates {
        (64.0,Ch 03: Lexical Scanning)
        (111.8,Ch 03: RegExp Tutorial)
        (229.1,Ch 03: Exam Evaluation)
        (56.2,Ch 04--05: NFA to DFA)
        (74.6,Ch 04--05: RegExp to NFA)
        (34.0,Ch 04--05: DFA to RegExp)
        (60.8,Ch 04--05: DFA Minimization)
        (489.4,Ch 04--05: Equivalence Check)
        (56.6,Ch 06: Top-Down Parser)
        (48.8,Ch 06: EBNF Parser)
        (87.1,Ch 07: Differentiator)
        (961.9,Ch 07/08: Interpreter)
        (108.3,Ch 08/09: Earley Parser)
        (647.6,Ch 10/11: Parser Conflicts)
        (1752.2,Ch 12/13: Integer-C Compiler)
      };
      \addplot+[fill=OCamlOrange!85, draw=OCamlOrange] coordinates {
        (390.6,Ch 03: Lexical Scanning)
        (959.2,Ch 03: RegExp Tutorial)
        (436.5,Ch 03: Exam Evaluation)
        (353.0,Ch 04--05: NFA to DFA)
        (389.8,Ch 04--05: RegExp to NFA)
        (323.7,Ch 04--05: DFA to RegExp)
        (368.0,Ch 04--05: DFA Minimization)
        (394.6,Ch 04--05: Equivalence Check)
        (463.2,Ch 06: Top-Down Parser)
        (454.4,Ch 06: EBNF Parser)
        (589.4,Ch 07: Differentiator)
        (1058.7,Ch 07/08: Interpreter)
        (533.5,Ch 08/09: Earley Parser)
        (832.0,Ch 10/11: Parser Conflicts)
        (1429.2,Ch 12/13: Integer-C Compiler)
      };
      \addplot+[fill=FSharpPurple!85, draw=FSharpPurple] coordinates {
        (1329.3,Ch 03: Lexical Scanning)
        (1495.2,Ch 03: RegExp Tutorial)
        (1575.7,Ch 03: Exam Evaluation)
        (1448.3,Ch 04--05: NFA to DFA)
        (1428.7,Ch 04--05: RegExp to NFA)
        (1372.6,Ch 04--05: DFA to RegExp)
        (1497.9,Ch 04--05: DFA Minimization)
        (1669.0,Ch 04--05: Equivalence Check)
        (1320.5,Ch 06: Top-Down Parser)
        (1293.8,Ch 06: EBNF Parser)
        (1567.5,Ch 07: Differentiator)
        (1881.0,Ch 07/08: Interpreter)
        (1634.3,Ch 08/09: Earley Parser)
        (1747.0,Ch 10/11: Parser Conflicts)
        (1892.1,Ch 12/13: Integer-C Compiler)
      };
      \legend{Python 3.13, OCaml 5.4, F\# 9.0}
    \end{axis}
  \end{tikzpicture}
  \caption{Comparison of cumulative code-cell execution times ($T_{\text{cells}}$) across 15 curriculum notebooks (excluding the SLR C-grammar outlier, analyzed separately in \cref{subsec:slr_case_study}).}
  \label{fig:macrobenchmarks_cells}
\end{figure}

\subsection{Why Compiled Languages Appear Slower on Short Notebooks---and Where They Overtake Python}

At first glance, \cref{tab:macrobenchmarks} presents an apparent paradox: although OCaml and F\# are statically typed, compiled languages whose native loops and ADT operations run $10\times$--$95\times$ faster than Python (\cref{sec:microbenchmarks}), Python achieves the lowest median notebook execution time ($T_{\text{wall}} = 599.2\text{ ms}$, $T_{\text{cells}} = 97.7\text{ ms}$), followed by OCaml ($T_{\text{wall}} = 1{,}169.4\text{ ms}$, $T_{\text{cells}} = 457.4\text{ ms}$) and F\# ($T_{\text{wall}} = 2{,}476.6\text{ ms}$, $T_{\text{cells}} = 1{,}532.7\text{ ms}$).

Decomposing the cell-by-cell execution profiles resolves this paradox completely:
\begin{enumerate}
  \item \textbf{In-Cell Preamble and Dependency Compilation Cost:}
  In every OCaml notebook, the opening cell executes \texttt{\#use "../lib/nb.ml";;}, which invokes the OCaml compiler to parse, type-check, and bytecode-compile \texttt{fset.ml}, \texttt{notebook.ml}, and \texttt{graphviz.ml} ($\approx 280\text{ ms}$). In every F\# notebook, the opening cell executes dependencies such as \texttt{FsLexYacc.Runtime} or project assemblies via \texttt{\#r} and \texttt{\#load}, triggering NuGet package resolution, assembly binding, and \texttt{FSharp.Compiler.Service} initialization ($\approx 1{,}150\text{ ms}$ in Cell~1).
  Once Cell~1 completes in F\#, the remaining 8 to 34 cells in a typical notebook execute in a combined $150\text{ ms}$--$350\text{ ms}$ ($\approx 10\text{ ms}$--$20\text{ ms}$ per cell, including type-checking and JIT compilation).

  \item \textbf{Crossover on Computationally Intensive Notebooks:}
  Whenever a notebook performs non-trivial algorithmic work or external compilation, OCaml and F\# rapidly close the gap or overtake Python outright:
  \begin{itemize}
    \item In \textbf{Ch 04--05: Equivalence Check} (\texttt{09-Equivalence.ipynb}), which constructs, determinizes, and minimizes full regular-expression automata to test language equivalence, Python's cell execution time jumps $8.7\times$ (from $56.2\text{ ms}$ to $489.4\text{ ms}$), whereas OCaml executes the entire notebook in $394.6\text{ ms}$ ($T_{\text{wall}} = 903.2\text{ ms}$ vs.\ Python's $1{,}002.0\text{ ms}$), beating Python in both cell execution and total wall-clock time!
    \item In \textbf{Ch 12/13: Integer-C Compiler} (\texttt{Compiler.ipynb}), which parses Integer-C programs, performs type checking, generates Jasmin JVM assembly, invokes the external Jasmin assembler, and executes the compiled \texttt{.class} files on the OpenJDK JVM, \textbf{OCaml achieves the fastest overall execution} ($T_{\text{cells}} = 1{,}429.2\text{ ms}$, $T_{\text{wall}} = 2{,}103.1\text{ ms}$), outperforming Python ($T_{\text{cells}} = 1{,}752.2\text{ ms}$, $T_{\text{wall}} = 2{,}303.4\text{ ms}$), while F\#'s cell time ($1{,}892.1\text{ ms}$) is nearly identical to Python despite F\#'s $1.1\text{ s}$ Cell~1 startup!
  \end{itemize}

  \item \textbf{Cold-Cache Sensitivity in Python Ecosystem Libraries:}
  On a newly provisioned virtual environment, Python's \texttt{03-Exam-Evaluation.ipynb} required $10{,}170.6\text{ ms}$ ($9{,}673.9\text{ ms}$ in one cell: \texttt{import matplotlib}) because \texttt{matplotlib} scans system fonts on macOS to build its cache. In contrast, neither OCaml nor F\# suffered library import stalls.
\end{enumerate}

\subsection{Case Study: The ANSI C \texorpdfstring{$\text{SLR}(1)$}{SLR(1)} Table Generation Stress Test}
\label{subsec:slr_case_study}

The single largest divergence in \cref{tab:macrobenchmarks} occurs in \textbf{Ch 09/10: SLR Table Generator},\\ where Python requires $1{,}058.9\text{ ms}$, OCaml requires $3{,}926.6\text{ ms}$, and F\# requires $14{,}842.0\text{ ms}$ (with peak RSS reaching $1{,}597.5\text{ MB}$).

Inspection of the notebook cells and underlying data structures reveals two compounding causes:
\begin{enumerate}
  \item \textbf{Doubled Workload in the F\# Verification Suite:}
  Whereas the Python and OCaml editions run the SLR(1) table generator once on the conflict-free ANSI~C grammar (\texttt{Examples/c-grammar-slr.g}, containing $65$ non-terminals and $220+$ productions), the F\# edition (\texttt{f\#/Chapter-09/SLR-Table-Generator.ipynb}, Cells 56--57) executes the full SLR(1) canonical item-set construction and action-table computation \emph{twice} on the entire C grammar: first on \texttt{Examples/c-grammar-slr.g} and a second time on the ambiguous \texttt{Examples/c-grammar.g} to assert that dangling-\texttt{else} shift/reduce conflicts are detected.

  \item \textbf{Structural Comparison of Deeply Nested Item Sets:}
  An SLR(1) state is a closure set of marked grammar rules (\texttt{Set<MarkedRule>}), and the canonical collection $\mathcal{C}$ is a set of states (\texttt{Set<Set<MarkedRule>>}), where each \texttt{MarkedRule} is a recursive discriminated union containing string symbols and lists of grammar symbols. For the ANSI~C grammar, $|\mathcal{C}| \approx 350$ states, each containing $20$--$150$ marked rules.
  As proven in our Workload~M2 analysis (\cref{sec:microbenchmarks}), Python represents each marked rule and state as a \texttt{frozenset} whose 64-bit hash is computed once and cached in C, making state deduplication $O(1)$ per closure step. In OCaml (\texttt{FSet}) and F\# (\texttt{Set<Set<MarkedRule>>}), every closure step inserts and compares sets inside balanced binary trees without cached hashes, performing millions of structural tree-to-tree comparisons over nested strings and lists.
\end{enumerate}

\begin{finding}[Algorithmic Impact of Structural Comparison on Nested Sets]
Representing canonical $\text{LR}/\text{SLR}$ state collections as un-hashed persistent trees incurs an $O(|I| \log |\mathcal{C}|)$ structural comparison penalty per state transition. Assigning integer state IDs or attaching a precomputed hash tag (hash-consing) to item sets in OCaml and F\# eliminates this bottleneck completely and restores native compiled speed.
\end{finding}

% ------------------------------------------------------------------------------
\section{Memory Residency and Garbage Collection Dynamics}
\label{sec:memory_gc}
% ------------------------------------------------------------------------------

Memory consumption is a critical operational metric both for personal student laptops (where running multiple Jupyter kernels concurrently can exhaust physical RAM) and for cloud-hosted JupyterHub deployments (which enforce strict per-container cgroup memory limits, typically $512\text{ MB}$--$2\text{ GB}$). \Cref{fig:memory_rss} contrasts the memory residency profiles of the three platforms across idle kernels, standalone algorithmic microbenchmarks, typical curriculum notebooks, and the heavy ANSI~C SLR(1) table generation workload.

\begin{figure}[htbp]
  \centering
  \begin{tikzpicture}
    \begin{axis}[
      ybar,
      ymode=log,
      log origin=infty,
      bar width=10pt,
      width=0.92\textwidth,
      height=7.2cm,
      ylabel={Resident Set Size (MB, log scale)},
      symbolic x coords={
        Microbench,
        Idle Kernel,
        Curriculum NB,
        Peak SLR
      },
      xtick=data,
      xticklabels={
        \thead{Standalone\\Microbench},
        \thead{Idle Jupyter\\Kernel},
        \thead{Median\\Curriculum NB},
        \thead{Peak\\(SLR C-Grammar)}
      },
      xticklabel style={font=\scriptsize, align=center},
      ymin=5,
      ymax=3000,
      ymajorgrids=true,
      grid style={dashed,gray!40},
      legend style={at={(0.03,0.95)}, anchor=north west, font=\footnotesize},
      nodes near coords,
      every node near coord/.append style={font=\tiny, rotate=0},
      enlarge x limits=0.18
    ]
      \addplot+[fill=PythonBlue!85, draw=PythonBlue] coordinates {
        (Microbench, 21.8)
        (Idle Kernel, 55.4)
        (Curriculum NB, 57.5)
        (Peak SLR, 68.5)
      };
      \addplot+[fill=OCamlOrange!85, draw=OCamlOrange] coordinates {
        (Microbench, 8.8)
        (Idle Kernel, 27.2)
        (Curriculum NB, 102.9)
        (Peak SLR, 242.0)
      };
      \addplot+[fill=FSharpPurple!85, draw=FSharpPurple] coordinates {
        (Microbench, 411.4)
        (Idle Kernel, 530.5)
        (Curriculum NB, 612.4)
        (Peak SLR, 1597.5)
      };
      \legend{Python 3.13, OCaml 5.4, F\# 9.0 (.NET 9)}
    \end{axis}
  \end{tikzpicture}
  \caption{Memory footprint (Resident Set Size in MB, log scale) across execution regimes for Python~3.13, OCaml~5.4, and F\#~9.0.}
  \label{fig:memory_rss}
\end{figure}

\subsection{Comparative Garbage Collector Architecture}

The memory profiles in \cref{fig:memory_rss} reflect the design philosophies of the three garbage collectors:

\begin{enumerate}
  \item \textbf{CPython 3.13 ($55.4\text{ MB}$ Idle $\to 57.5\text{ MB}$ Median $\to 105.5\text{ MB}$ Max):}
  CPython employs immediate reference counting (\texttt{Py\_INCREF} / \texttt{Py\_DECREF}) as its primary reclamation mechanism, supplemented by a three-generation cyclic garbage collector for container objects. Because reference counts drop to zero immediately when temporary sets or closures go out of scope at the end of a function or cell, memory blocks are returned to CPython's \texttt{pymalloc} object arenas instantaneously without waiting for a tracing GC sweep. As a result, Python's RSS remains remarkably flat across 13 of the 16 notebooks ($55.5\text{ MB}$--$68.5\text{ MB}$), rising to $91$--$105\text{ MB}$ only when loading \texttt{Lark}, \texttt{PLY}, or \texttt{matplotlib}.

  \item \textbf{OCaml 5.4 ($8.8\text{ MB}$ Standalone $\to 27.2\text{ MB}$ Idle $\to 102.9\text{ MB}$ Median):}
  OCaml~5 utilizes a generational copying/mark-and-sweep collector optimized for high-rate functional allocation~\cite{leroy2023ocaml,sivaramakrishnan2020retrofitting}. New variant values and AVL tree nodes are allocated on a fixed-size per-domain \emph{minor heap} via a single register increment; most short-lived intermediate AST and set nodes die before the minor heap fills, incurring zero scavenging cost. In a standalone binary, OCaml requires just $8.77\text{ MB}$ of RAM. Inside \texttt{ocaml-jupyter}, loading the compiler libraries (\texttt{compiler-libs.toplevel}), \texttt{merlin} type tables, and PCRE2/Yojson brings the steady-state notebook footprint to a highly predictable $102.2\text{ MB}$--$107.4\text{ MB}$.

  \item \textbf{F\# 9.0 on .NET 9 CoreCLR ($530.5\text{ MB}$ Idle $\to 612.4\text{ MB}$ Median $\to 1{,}597.5\text{ MB}$ Peak):}
  The .NET~9 CoreCLR runtime paired with \texttt{dotnet-interactive} consumes roughly $6\times$--$10\times$ more resident memory than Python and OCaml. This footprint is driven by three factors:
  \begin{itemize}
    \item \emph{In-Memory Compiler Service Residency:} Unlike Python (which discards ASTs after bytecode generation) and OCaml (whose top-level keeps compact \texttt{.cmi} module signatures), \texttt{FSharp.Compiler.Service} retains full typed abstract syntax trees (TASTs), assembly metadata readers, and cross-cell incremental type-checking caches in memory for IDE/LSP intellisense.
    \item \emph{Uncollectible Dynamic Assemblies:} Every executed cell in \texttt{dotnet-interactive} is compiled into a distinct dynamic CLI class/module loaded into the current \texttt{AppDomain} / \texttt{AssemblyLoadContext}, accumulating JIT-compiled native code pages and reflection metadata across the notebook session.
    \item \emph{Throughput-Oriented Workstation GC:} By default on modern multi-core hardware with ample physical RAM, the CoreCLR Garbage Collector delays full Gen~2 compacting collections in favor of high allocation throughput, allowing the ephemeral and Gen~1/Gen~2 heaps to expand to $1.6\text{ GB}$ during the heavy ANSI~C SLR(1) closure computation rather than throttling execution with frequent stop-the-world collections.
  \end{itemize}
\end{enumerate}

% ------------------------------------------------------------------------------
\section{Static Corpus Metrics and Execution Correlation}
\label{sec:corpus_correlation}
% ------------------------------------------------------------------------------

To contextualize how code structure interacts with notebook execution times, \cref{tab:corpus_metrics} summarizes the static structural metrics across the entire repository (\texttt{code\_metrics.py}).

\begin{table}[htbp]
  \centering
  \caption{Static corpus inventory across the Python, OCaml, and F\# notebook editions.}
  \label{tab:corpus_metrics}
  \small
  \setlength{\tabcolsep}{5pt}
  \begin{tabular}{@{}lrrrrrr@{}}
    \toprule
    \thead{Edition} &
    \thead{Notebooks} &
    \thead{Code\\Cells} &
    \thead{Markdown\\Cells} &
    \thead{Total\\SLOC} &
    \thead{SLOC /\\Notebook} &
    \thead{Cells /\\Notebook} \\
    \midrule
    Python 3.13 & $49$ & $1{,}013$ & $936$ & $5{,}333$ & $108.8$ & $20.7$ \\
    OCaml 5.4   & $50$ & $991$     & $968$ & $6{,}453$ & $129.1$ & $19.8$ \\
    F\# 9.0     & $56$ & $667$     & $667$ & $3{,}486$ & $62.2$  & $11.9$ \\
    \bottomrule
  \end{tabular}
\end{table}

As shown in \cref{tab:corpus_metrics}, the F\# edition achieves a $34.6\%$ reduction in total notebook SLOC ($3{,}486$ vs.\ $5{,}333$ in Python and $6{,}453$ in OCaml) and averages $11.9$ code cells per notebook compared to $\approx 20$ in Python and OCaml. This structural compression arises from two architectural choices in the F\# port:
\begin{enumerate}
  \item Concise algebraic data types and pipeline expressions (\texttt{Set.fold}, \texttt{Map.tryFind}, \texttt{|>}) replace multi-line imperative loops; and
  \item Lexer and parser specifications (\texttt{.fsl} and \texttt{.fsy} files) are compiled into assemblies (\texttt{Scanning.dll}, \texttt{CalculatorSyntax.dll}, etc.) via \texttt{dotnet build} and referenced via \texttt{\#r}, whereas Python and OCaml embed grammar strings inline and invoke \texttt{Lark}/\texttt{PLY} or \texttt{ocamllex}/\texttt{Menhir} during cell execution.
\end{enumerate}

% ------------------------------------------------------------------------------
\section{Synthesis of Empirical Performance Findings}
\label{sec:perf_synthesis}
% ------------------------------------------------------------------------------

\Cref{tab:performance_synthesis} synthesizes our empirical findings across all evaluated runtime dimensions.

\begin{table}[htbp]
  \centering
  \caption{Comparative performance synthesis across runtime regimes.}
  \label{tab:performance_synthesis}
  \scriptsize
  \setlength{\tabcolsep}{3pt}
  \begin{tabularx}{\textwidth}{@{}llllX@{}}
    \toprule
    \thead{Performance Dimension} &
    \thead{Python 3.13} &
    \thead{OCaml 5.4} &
    \thead{F\# 9.0} &
    \thead{Architectural Root Cause} \\
    \midrule
    \textbf{Kernel Cold Start} &
      \textcolor{ForestGreen}{\textbf{Fast}} ($396\text{ ms}$) &
      \textcolor{ForestGreen}{\textbf{Fast}} ($511\text{ ms}$) &
      \textcolor{BrickRed}{Slow} ($1{,}884\text{ ms}$) &
      FCS symbol loading \& .NET BCL metadata resolution \\
    \addlinespace[3pt]
    \textbf{Warm Cell IPC Floor} &
      $4.49\text{ ms}$ &
      \textcolor{ForestGreen}{\textbf{Best}} ($3.50\text{ ms}$) &
      \textcolor{ForestGreen}{\textbf{Fast}} ($3.70\text{ ms}$) &
      Native ZeroMQ handling after JIT warmup \\
    \addlinespace[3pt]
    \textbf{Lexical Scanning (M1)} &
      \textcolor{BrickRed}{Slow} ($62.1\text{ ms}$) &
      \textcolor{ForestGreen}{\textbf{Fast}} ($3.5\text{ ms}$) &
      \textcolor{ForestGreen}{\textbf{Best}} ($1.9\text{ ms}$) &
      RyuJIT loop optimization vs.\ \texttt{re.Match} boxing \\
    \addlinespace[3pt]
    \textbf{Powerset $\mathcal{P}(Q)$ (M2)} &
      \textcolor{ForestGreen}{\textbf{Best}} ($9.1\text{ ms}$) &
      \textcolor{ForestGreen}{\textbf{Fast}} ($10.7\text{ ms}$) &
      \textcolor{BrickRed}{Slow} ($58.5\text{ ms}$) &
      $O(1)$ cached hash on \texttt{frozenset} vs.\ $O(k^2)$ AVL tree compare \\
    \addlinespace[3pt]
    \textbf{AST Rewriting (M4)} &
      \textcolor{BrickRed}{Slow} ($1.42\text{ ms}$) &
      \textcolor{ForestGreen}{\textbf{Best}} ($0.015\text{ ms}$) &
      \textcolor{ForestGreen}{\textbf{Fast}} ($0.11\text{ ms}$) &
      Minor-heap bump allocation \& compiled jump-table matching \\
    \addlinespace[3pt]
    \textbf{Short Notebook Wall Time} &
      \textcolor{ForestGreen}{\textbf{Best}} ($0.60\text{ s}$) &
      Moderate ($1.17\text{ s}$) &
      Higher ($2.48\text{ s}$) &
      Dominated by fixed Cell~1 compiler/NuGet initialization \\
    \addlinespace[3pt]
    \textbf{Heavy Compiler Notebook} &
      $2.30\text{ s}$ &
      \textcolor{ForestGreen}{\textbf{Best}} ($2.10\text{ s}$) &
      $2.91\text{ s}$ &
      OCaml native/bytecode speed overtakes Python on large tasks \\
    \addlinespace[3pt]
    \textbf{Memory Footprint (RSS)} &
      \textcolor{ForestGreen}{\textbf{Low}} ($58\text{ MB}$) &
      \textcolor{ForestGreen}{\textbf{Low}} ($103\text{ MB}$) &
      \textcolor{BrickRed}{High} ($612\text{ MB}$) &
      Reference counting / compact heap vs.\ hosted Compiler Service \\
    \bottomrule
  \end{tabularx}
\end{table}

In conclusion, our empirical measurements demonstrate that \textbf{no single language dominates across all performance regimes}:
\begin{itemize}
  \item \textbf{OCaml~5.4} offers the most balanced raw execution profile: near-instantaneous AST manipulation ($95\times$ faster than Python), low and predictable memory residency ($27\text{ MB}$ idle, $103\text{ MB}$ active), the fastest warm cell latency ($3.50\text{ ms}$), and the fastest end-to-end execution on complex multi-stage compiler notebooks ($2.10\text{ s}$).
  \item \textbf{Python~3.13} excels in rapid kernel startup ($396\text{ ms}$), low memory overhead ($58\text{ MB}$ median RSS), and set-of-sets fixed-point constructions where C-backed \texttt{frozenset} hash caching provides an unexpected asymptotic constant-factor advantage over purely functional balanced binary trees.
  \item \textbf{F\#~9.0 (.NET~9)} delivers the highest raw loop and scanner throughput ($1.92\text{ ms}$ on $300{,}000$ tokens via RyuJIT) and competitive warm cell execution, but pays a fixed startup ($1.88\text{ s}$) and memory ($530$--$612\text{ MB}$) penalty inside \texttt{dotnet-interactive} due to in-process compiler hosting, and requires custom hash-consing when operating on deeply nested \texttt{Set<Set<'T>>} structures.
\end{itemize}
